\documentclass[preprint,12pt]{elsarticle}

\usepackage{booktabs}
\usepackage{multirow}
\usepackage{graphicx}
\usepackage{xcolor}
\usepackage{hyperref}
\usepackage{amsmath}
\usepackage{enumitem}
\usepackage{tabularx}
\usepackage{siunitx}
\usepackage{url}

\hypersetup{colorlinks=true, linkcolor=blue, citecolor=blue, urlcolor=blue}

\journal{Journal of Systems and Software}

\begin{document}

\begin{frontmatter}

\title{Benchmarking JavaScript Bundlers at Scale:\\A Two-Tier Empirical Study of Vite, Rspack, esbuild, Webpack, and Rollup}

\author[addr1]{[Author Name]}
\address[addr1]{[Affiliation]}

\begin{abstract}
JavaScript bundlers are a critical component of modern web development, transforming source modules into optimized production artifacts.
Despite their ubiquity---Staicu et al.\ found that 40\% of websites rely on bundlers---no peer-reviewed empirical comparison of modern bundlers exists.
Existing benchmarks suffer from trivially synthetic projects, vendor maintenance bias, and absence of statistical rigor.

We present the first rigorous empirical comparison of five production-ready JavaScript bundlers---Vite~8, Rspack~2, esbuild~0.28, Webpack~5, and Rollup~4---across 11 directly measured metrics spanning build performance, output quality, and resource consumption.
Our study employs a two-tier benchmarking strategy: Tier~1 uses calibrated synthetic React+TypeScript projects at five scales (50--5{,}000 modules) for controlled comparison; Tier~2 validates findings on a real-world open-source project (Bulletproof React, 102 components).
We apply non-parametric statistics---Kruskal--Wallis omnibus tests, Mann--Whitney $U$ pairwise comparisons with Bonferroni correction, and Cliff's delta effect sizes---to 1{,}534 valid data points.

Results show that tool selection depends on project scale: esbuild is fastest at $\leq$500 modules, while Vite~8 (powered by its Rust-based Rolldown engine) dominates at $\geq$2{,}000 modules, completing a 5{,}000-module production build in 5.9\,s versus 73.5\,s for Webpack.
All five tools exhibit linear build-time scaling ($R^2 \geq 0.97$), with Vite's slope 23.9$\times$ lower than Webpack's.
Of 27 omnibus tests, all are statistically significant ($p < 0.05$); 197 of 205 pairwise comparisons yield large effect sizes.
Tier~2 rankings match Tier~1, supporting external generalizability.
We provide a full replication package.
\end{abstract}

\begin{keyword}
JavaScript bundler \sep build tools \sep empirical study \sep performance benchmarking \sep Vite \sep Rspack \sep esbuild
\end{keyword}

\begin{highlights}
\item First rigorous empirical comparison of five JavaScript bundlers
\item Two-tier design: calibrated synthetic + real-world OSS validation
\item 1{,}534 measurements across 11 metrics and 5 project scales
\item Vite 8 overtakes esbuild between 500 and 2{,}000 modules
\item Full replication package with scripts, data, and analysis code
\end{highlights}

\end{frontmatter}

%% ═══════════════════════════════════════════════════════════════
%% SECTION 1: INTRODUCTION (placeholder)
%% ═══════════════════════════════════════════════════════════════
\section{Introduction}\label{sec:intro}

\textit{[To be written --- see WRITING-PLAN.md \S1]}

%% ═══════════════════════════════════════════════════════════════
%% SECTION 2: BACKGROUND (placeholder)
%% ═══════════════════════════════════════════════════════════════
\section{Background}\label{sec:background}

\textit{[To be written --- see WRITING-PLAN.md \S2]}

%% ═══════════════════════════════════════════════════════════════
%% SECTION 3: RELATED WORK (placeholder)
%% ═══════════════════════════════════════════════════════════════
\section{Related Work}\label{sec:related}

\textit{[To be written --- see WRITING-PLAN.md \S3]}


%% ═══════════════════════════════════════════════════════════════
%% SECTION 4: STUDY DESIGN
%% ═══════════════════════════════════════════════════════════════
\section{Study Design}\label{sec:design}

This section describes the design of our empirical study, following the guidelines of Wohlin et al.~\cite{wohlin2012experimentation} and the benchmarking methodology established by Kalibera and Jones~\cite{kalibera2013rigorous}.

\subsection{Research Questions}\label{sec:rqs}

We structure our investigation around three research questions, each addressing a distinct dimension of bundler performance.

\medskip\noindent\textbf{RQ1: How do modern JavaScript bundlers compare in build performance across project scales?}
Build speed directly impacts developer productivity and is the primary reason teams migrate bundlers~\cite{stateofjs2025}.
We measure four metrics: development server cold-start time (M1), production build time (M2), incremental rebuild latency (M3), and hot module replacement (HMR) latency (M4).

\medskip\noindent\textbf{RQ2: How does bundler choice affect output quality and delivery efficiency?}
A fast bundler that produces bloated output is a false economy; bundle size and code-splitting granularity directly affect end-user load times.
We measure raw bundle size (M5), gzipped bundle size (M6), tree-shaking effectiveness (M7), code-splitting granularity (M8), and sourcemap accuracy (M9).

\medskip\noindent\textbf{RQ3: How do resource consumption and scaling behavior differ across bundlers?}
Resource usage determines viability in CI/CD environments with memory limits and on developer machines with 8--32\,GB RAM.
We measure peak resident-set-size memory (M10), CPU time (M11), and derive scaling regression models from M2 across all five project scales (M12).

\subsection{Tool Selection}\label{sec:tools}

We select five tools spanning three architectural paradigms and three implementation languages, ensuring a balanced comparison.
The inclusion criteria are: (1)~open source and installable via npm; (2)~standalone (usable outside a specific framework); (3)~production-ready with stable release; (4)~significant adoption ($\geq$100K weekly npm downloads or major organizational backing); and (5)~architecturally distinct from other selected tools.

Table~\ref{tab:tools} summarizes the selected tools.

\begin{table}[htbp]
\centering
\caption{Tools under study. All versions pinned on experiment Day~1 (September~11, 2026). Downloads from npm trends~\protect\cite{npmtrends2026}.}\label{tab:tools}
\begin{tabular}{@{}llllr@{}}
\toprule
\textbf{Tool} & \textbf{Version} & \textbf{Lang.} & \textbf{Paradigm} & \textbf{npm wk.\ DL} \\
\midrule
Vite     & 8.3.0   & Rust (Rolldown) & Unbundled dev  & $\sim$16M \\
Rspack   & 2.2.3   & Rust            & Fully bundled  & $\sim$700K \\
esbuild  & 0.28.2  & Go              & Fully bundled  & $\sim$38M \\
Webpack  & 5.110.3 & JS + SWC (Rust) & Fully bundled  & $\sim$28M \\
Rollup   & 4.63.1  & JavaScript      & Fully bundled  & $\sim$30M \\
\bottomrule
\end{tabular}
\end{table}

\textbf{Vite~8}~\cite{vite2026} represents the dominant greenfield choice ($\sim$16M weekly downloads). Version~8 is architecturally significant: it replaced both esbuild (development transforms) and Rollup (production bundling) with Rolldown~\cite{rolldown2026}, a unified Rust-based engine---a fundamental shift worth quantifying.

\textbf{Rspack~2}~\cite{rspack2026}, developed by ByteDance, offers a drop-in Webpack replacement with a Rust core. It represents the ``migrate existing Webpack projects to Rust'' migration path and tests whether API compatibility incurs a performance cost.

\textbf{esbuild~0.28}~\cite{esbuild2026}, written in Go, pioneered the 10--100$\times$ speedup narrative that catalyzed the Rust/Go bundler wave. Its minimal plugin API tests the raw-speed-versus-extensibility tradeoff. As the only Go-based tool, it provides a cross-language comparison point.

\textbf{Webpack~5}~\cite{webpack2026} is the incumbent, used by $\sim$86\% of developers~\cite{stateofjs2025}. We equip it with SWC~\cite{swc2026} for Rust-accelerated transforms, giving it the best-case scenario. It serves as the \emph{baseline control}.

\textbf{Rollup~4}~\cite{rollup2026}, Vite's former production engine, is a pure-JavaScript ESM-native bundler. It serves as the \emph{second baseline}, establishing the performance floor of a tool without Rust or Go acceleration.

\medskip\noindent\textbf{Excluded tools.} Turbopack lacks a stable standalone production build (it functions only as a Next.js-internal development bundler) and is noted as future work. Parcel~2 overlaps architecturally with Vite. Farm has insufficient adoption ($\sim$2K GitHub stars). Bun's bundler is runtime-integrated (different execution model). SWC is a transpiler, not a full bundler.

\subsection{Two-Tier Benchmarking Strategy}\label{sec:twotier}

A key design decision is how to construct the projects being benchmarked.
Existing vendor-maintained benchmarks (e.g., rspack-contrib/build-tools-performance~\cite{rspackbench2026}, Rolldown benchmarks~\cite{rolldownbench2026}) use trivially synthetic components ($\sim$5 lines of empty JSX), while purely real-world projects cannot be fairly ported to all five bundlers due to bundler-specific APIs.
Our solution is a \textbf{two-tier strategy} that satisfies both internal validity (controlled comparison via synthetic projects) and external validity (real-world generalizability).

\subsubsection{Tier~1: Calibrated Synthetic Projects}\label{sec:tier1}

We developed a deterministic project generator\footnote{Included in the replication package; seeded with Mulberry32 PRNG (seed = 42) for exact reproducibility.} that produces React + TypeScript projects at five scales, summarized in Table~\ref{tab:projects}.

\begin{table}[htbp]
\centering
\caption{Tier~1 synthetic project specifications. All projects generated with seed = 42. Dead modules (imported by nothing) are fixed at 5\% to exercise tree-shaking. LOC counts include both TSX component code and CSS module files.}\label{tab:projects}
\begin{tabular}{@{}lrrrrrrr@{}}
\toprule
\textbf{Size} & \textbf{Modules} & \textbf{Folders} & \textbf{Routes} & \textbf{Import Edges} & \textbf{Dead (\%)} & \textbf{LOC} \\
\midrule
xs   &     50 &  5 &   0 &     80 &  2 (4.0\%) &    8{,}135 \\
s    &    200 & 10 &   4 &    369 & 10 (5.0\%) &   32{,}209 \\
m    &    500 & 15 &  12 &    898 & 25 (5.0\%) &   79{,}605 \\
l    &  2{,}000 & 30 &  20 &  3{,}647 & 100 (5.0\%) &  319{,}647 \\
xl   &  5{,}000 & 50 &  50 &  9{,}301 & 250 (5.0\%) &  801{,}736 \\
\bottomrule
\end{tabular}
\end{table}

Each generated component falls into one of eight template categories---data display (25\%), form control (15\%), layout (15\%), navigation (10\%), feedback (10\%), data fetching (10\%), chart/visualization (10\%), and utility/context providers (5\%)---weighted to approximate the distribution in production codebases.
Components range from 50 to 150 lines of code and include realistic patterns: React hooks (\texttt{useState}, \texttt{useMemo}, \texttt{useCallback}, \texttt{useEffect}), third-party imports (\texttt{lodash-es}, \texttt{date-fns}, \texttt{clsx}), CSS Modules, TypeScript generics, and inter-component dependencies.
Critically, each component template is \emph{calibrated}---designed so that specific features exercise specific metrics.
For example, the \texttt{lodash-es} import (using only \texttt{pick} and \texttt{debounce}) exercises tree-shaking (M7), lazy-loaded route entries exercise code splitting (M8), and dead modules (5\% of components, imported by nothing) test dead-code elimination.

The dependency graph is structured to resemble real projects: each component imports 1--3 siblings within its feature folder; 10\% of components have a cross-folder import; and $\sim$2\% are designated as ``hub'' components imported by 20--30 parents.
Circular dependencies are detected and broken by the generator.

\subsubsection{Tier~2: Real-World OSS Validation}\label{sec:tier2}

To validate that synthetic results generalize, we select Bulletproof React~\cite{bulletproofreact2024}, an open-source production-architecture reference application with 102 TSX component files, React 18, TypeScript 5, Tailwind CSS, React Query, Zustand, and React Router.
We adapt it for all five bundlers by creating tool-specific configurations that follow each bundler's official documentation, with no tool-specific optimizations beyond defaults.
All configurations are included in the replication package.

\subsection{Metrics}\label{sec:metrics}

Table~\ref{tab:metrics} defines the 11 directly measured metrics (M1--M11) and one derived metric (M12).
Metrics M1--M4 capture build performance (RQ1), M5--M9 capture output quality (RQ2), and M10--M12 capture resource consumption and scaling (RQ3).

\begin{table}[htbp]
\centering
\caption{Metric definitions. Run counts indicate the number of independent measurements per tool$\times$size configuration.
M5--M9 are single-run deterministic metrics (verified via a second build).
M12 is derived from M2 medians across all five Tier~1 scales.}\label{tab:metrics}
\begin{tabular}{@{}cllllr@{}}
\toprule
\textbf{ID} & \textbf{Name} & \textbf{Unit} & \textbf{Measurement Tool} & \textbf{Scope} & \textbf{Runs} \\
\midrule
M1  & Dev cold start     & ms      & Puppeteer + log parsing & V, R$^*$, W  & 20 \\
M2  & Prod build time    & ms      & hyperfine               & All 5        & 10 \\
M3  & Incremental rebuild & ms     & \texttt{build --watch} + fs event  & V, R, W, Ro  & 20 \\
M4  & HMR latency        & ms      & Puppeteer + CDP console & V, R, W      & 20 \\
\midrule
M5  & Bundle size (raw)  & bytes   & \texttt{find dist | wc -c}  & All 5    & 1 \\
M6  & Bundle size (gzip) & bytes   & \texttt{gzip -9 -c | wc -c} & All 5   & 1 \\
M7  & Tree-shaking        & count  & grep lodash references   & All 5       & 1 \\
M8  & Code splitting     & count   & JS chunk count in dist   & All 5       & 1 \\
M9  & Sourcemap accuracy & ratio   & map count / JS count     & All 5       & 1 \\
\midrule
M10 & Peak memory (RSS)  & MB      & \texttt{command time -l}  & All 5       & 10 \\
M11 & CPU time           & s       & user + sys from \texttt{time} & All 5  & 10 \\
M12 & Scaling regression & ---     & Derived: $R^2$, slope    & All 5       & --- \\
\bottomrule
\multicolumn{6}{@{}l@{}}{\footnotesize V = Vite, R = Rspack, W = Webpack, Ro = Rollup. $^*$Rspack M1 timed out at all sizes.}
\end{tabular}
\end{table}

\subsection{Measurement Harness}\label{sec:harness}

All experiments were conducted on a single machine to eliminate hardware variability: Apple M2~Pro (12-core, 32\,GB RAM) running macOS~26.5, Node.js~v22.16.0, and npm~10.9.2.
Background processes were minimized during all measurement runs.

\textbf{M2 (production build time)} is measured using hyperfine~\cite{hyperfine2024}, a command-line benchmarking tool that reports per-run wall-clock times.
Before each run, a cache-clearing script removes \texttt{node\_modules/.cache}, \texttt{dist/}, \texttt{.rspack/}, \texttt{.vite/}, and OS filesystem caches where applicable.
We configure hyperfine with zero explicit warm-up runs (\texttt{--warmup 0}) and 10 measured runs (\texttt{--runs 10}).

\textbf{M1 (dev cold start)} is measured by starting the development server process, monitoring \texttt{stdout} for the tool-specific ``ready'' pattern (e.g., \texttt{"Local:"} for Vite, \texttt{"compiled"} for Webpack/Rspack), and recording the elapsed wall-clock time.
This is repeated for 20 independent runs with cache clearing between each.
A 60-second timeout classifies non-responding servers as failures.

\textbf{M3 (incremental rebuild)} uses each tool's \texttt{build --watch} mode.
A dedicated Node.js script starts the watch process, waits for the initial build to complete, then programmatically modifies a randomly selected component file (toggling a JSX attribute) and measures the time until the watcher reports rebuild completion.
This is repeated for 20 file changes per configuration.

\textbf{M4 (HMR latency)} is measured via Puppeteer~\cite{puppeteer2024} controlling a headless Chromium browser.
The script launches the development server, navigates Puppeteer to the local URL, intercepts Chrome DevTools Protocol (CDP) console messages, then modifies a component file.
The HMR latency is the interval between the file-system write and the browser console's HMR confirmation message.
This is repeated for 20 independent file modifications.

\textbf{M5--M9 (output quality)} are deterministic single-run metrics computed from the production build output.
Each is verified by a second independent build to confirm determinism.
M5 sums all file sizes in the output directory; M6 applies \texttt{gzip -9} to each JavaScript file and sums; M7 counts residual \texttt{lodash-es} references (full library is $\sim$600\,KB; our code imports only \texttt{pick} and \texttt{debounce}); M8 counts distinct JavaScript chunks; M9 reports the ratio of sourcemap files to JavaScript files.

\textbf{M10 and M11 (memory and CPU)} are extracted from the macOS \texttt{command time -l} utility, which reports peak resident set size (in bytes) and user + system CPU time.
Ten independent runs are performed with cache clearing between each.

\textbf{No explicit warm-up.} Following Barrett et al.~\cite{barrett2017vm}, who showed that 43.5\% of VM+benchmark pairs fail to reach steady state, we do not assume warm-up convergence. Instead, we use the median (robust to outliers) as our primary summary statistic and report the full interquartile range.

\subsection{Statistical Analysis}\label{sec:stats}

Our statistical pipeline follows the recommendations of Georges et al.~\cite{georges2007statistically} for non-parametric analysis of performance data.

\textbf{Step 1: Normality testing.}
We apply the Shapiro--Wilk test~\cite{shapiro1965analysis} ($\alpha = 0.05$) to every tool$\times$size$\times$metric group with $n \geq 3$.
Of the 115 testable groups, 74 (64\%) reject the null hypothesis of normality, motivating non-parametric methods throughout.

\textbf{Step 2: Omnibus test.}
For each metric$\times$size combination where multiple tools have data, we apply the Kruskal--Wallis $H$ test~\cite{kruskal1952use} ($\alpha = 0.05$) to determine whether at least one tool differs significantly from the others.

\textbf{Step 3: Pairwise comparisons.}
Where the omnibus test is significant, we conduct all $\binom{k}{2}$ pairwise Mann--Whitney $U$ tests~\cite{mann1947test} with Bonferroni correction~\cite{bonferroni1936teoria} to control the family-wise error rate.

\textbf{Step 4: Effect sizes.}
For every pairwise comparison, we compute Cliff's delta~\cite{cliff1993dominance}, a non-parametric effect-size measure interpreted as: $|\delta| < 0.147$ (negligible), $< 0.33$ (small), $< 0.474$ (medium), and $\geq 0.474$ (large), following the thresholds of Vargha and Delaney~\cite{vargha2000critique}.

\textbf{Step 5: Scaling regression.}
For each tool, we fit both a linear model ($\text{build\_time} = a \cdot \text{modules} + b$) and a log-linear model ($\text{build\_time} = a \cdot \ln(\text{modules}) + b$) to the five Tier~1 M2 medians.
We select the model with the higher $R^2$.

All analysis code is implemented in TypeScript (executed via \texttt{tsx}) using the \texttt{simple-statistics} library and is included in the replication package.


%% ═══════════════════════════════════════════════════════════════
%% SECTION 5: RESULTS
%% ═══════════════════════════════════════════════════════════════
\section{Results}\label{sec:results}

This section presents results organized by research question.
All 27 Kruskal--Wallis omnibus tests are statistically significant ($p < 0.005$), confirming that at least one tool differs from the others in every metric$\times$size combination tested.
Of 205 pairwise Mann--Whitney $U$ comparisons with Bonferroni correction, 188 (91.7\%) are statistically significant.
Of 205 Cliff's delta effect sizes, 197 (96.1\%) are large, 4 are medium, 2 are small, and 2 are negligible.
All non-significant pairs involve adjacent-ranked tools whose medians are within practical noise.

\subsection{RQ1: Build Performance}\label{sec:rq1}

\subsubsection{M2: Production Build Time}

Table~\ref{tab:m2} presents the primary build-time results.

\begin{table}[htbp]
\centering
\caption{M2: Production build time in milliseconds (median [IQR], $n = 10$ runs per cell). Bold indicates the fastest tool at each scale.}\label{tab:m2}
\begin{tabular}{@{}lrrrrr@{}}
\toprule
\textbf{Tool} & \textbf{xs (50)} & \textbf{s (200)} & \textbf{m (500)} & \textbf{l (2{,}000)} & \textbf{xl (5{,}000)} \\
\midrule
esbuild  & \textbf{399} [41]  & \textbf{544} [56]  & \textbf{988} [35] & 3{,}592 [714]   & 14{,}682 [3{,}925] \\
Vite     & 1{,}017 [38]       & 1{,}431 [476]      & 1{,}644 [49]      & \textbf{2{,}926} [89] & \textbf{5{,}900} [379] \\
Rspack   & 1{,}150 [132]      & 1{,}518 [113]      & 2{,}516 [177]     & 6{,}134 [193]   & 13{,}848 [499] \\
Rollup   & 2{,}858 [37]       & 3{,}860 [73]       & 7{,}670 [191]     & 24{,}608 [676]  & 99{,}208 [1{,}202] \\
Webpack  & 3{,}712 [156]      & 5{,}214 [146]      & 10{,}982 [485]    & 39{,}664 [2{,}301] & 115{,}872 [8{,}416] \\
\bottomrule
\end{tabular}
\end{table}

\textbf{The dominant result is a scale-dependent crossover}: esbuild is the fastest tool at $\leq$500 modules, completing a 500-module build in 988\,ms---40\% faster than Vite's 1{,}644\,ms.
However, at $\geq$2{,}000 modules, Vite overtakes esbuild: at the xl-5{,}000 scale, Vite builds in 5.9\,s versus esbuild's 14.7\,s, a 2.5$\times$ advantage.
Linear interpolation of medians places the crossover at approximately 1{,}250 modules, though our data can only definitively confirm it occurs between 500 and 2{,}000 modules.

Rspack occupies a consistent third place, 1.5--2.3$\times$ slower than the leader at each scale.
Notably, the Rspack--Vite difference is not statistically significant at sizes s-200 ($p_\text{adj} = 1.0$) and xs-50 ($p_\text{adj} = 0.156$), with negligible-to-small Cliff's delta values, indicating that these two Rust-based tools perform comparably at small scales.

The JavaScript-based tools (Rollup and Webpack) are 2.9--19.7$\times$ slower than the fastest tool.
At xl-5{,}000, Webpack requires 115.9\,s and Rollup 99.2\,s.
The Rollup--Webpack difference at xl-5{,}000 is not statistically significant ($p_\text{adj} = 0.156$), suggesting both hit similar scalability ceilings at enterprise scale.

\subsubsection{M1: Development Server Cold Start}

M1 is limited to Vite and Webpack; Rspack timed out (60\,s) at all three tested sizes, which is itself a finding---its development server startup is prohibitively slow for the configurations used.

Vite's dev cold start is remarkably \emph{scale-invariant}: 868\,ms at xs-50, 870\,ms at m-500, and 863\,ms at xl-5{,}000 (IQR $\leq$ 31\,ms).
This reflects Vite's unbundled development architecture, which serves native ES modules on demand rather than pre-bundling the entire project.
In contrast, Webpack scales with project size: 2{,}485\,ms at xs-50, 4{,}513\,ms at m-500, and 18{,}924\,ms at xl-5{,}000---a 7.6$\times$ increase.
All M1 comparisons between Vite and Webpack are significant ($p < 10^{-6}$) with large effect sizes ($|\delta| = 1.0$).

\subsubsection{M3: Incremental Rebuild}

In watch mode, Webpack is the fastest at the smallest scale (138\,ms at xs-50), with Rspack close behind (188\,ms).
At m-500, Rspack (328\,ms) and Webpack (525\,ms) lead, while Vite (872\,ms) and Rollup (3{,}130\,ms) are slower.
The Rspack--Webpack difference at m-500 is marginally non-significant ($p_\text{adj} = 0.077$; medium $|\delta| = 0.46$).
esbuild is excluded from M3 as it lacks a native watch mode.
At xl-5{,}000, only Vite reported a single successful rebuild (16{,}547\,ms); all other tools either timed out or failed to detect the incremental change, indicating an industry-wide gap in watch-mode scalability at enterprise scale.

\subsubsection{M4: HMR Latency}

Rspack and Webpack achieve near-identical HMR latencies: 14\,ms at both xs-50 and m-500 for both tools, with negligible-to-small Cliff's delta values between them ($|\delta| = 0.34$ at xs-50, $|\delta| = 0.15$ at m-500).
Vite's HMR is slower at 76\,ms (m-500) and 132\,ms (xs-50, $n = 3$).
This is unexpected given Vite's native ESM architecture; the overhead likely stems from Vite's HMR protocol requiring a full module graph traversal, whereas Rspack and Webpack use direct chunk replacement.

\medskip\noindent\textbf{Finding 1.} \textit{Build tool selection for speed depends on project scale. esbuild dominates at $\leq$500 modules; Vite dominates at $\geq$2{,}000 modules. For incremental development, Rspack and Webpack offer the fastest watch-mode rebuilds and HMR, while Vite's dev server cold-start is uniquely scale-invariant.}


\subsection{RQ2: Output Quality}\label{sec:rq2}

Table~\ref{tab:quality} presents M5--M9 results.
These are deterministic single-run metrics, verified by a second independent build.

\begin{table}[htbp]
\centering
\caption{Output quality metrics at the m-500 and xl-5{,}000 scales (representative). M5/M6 in KB; M7 = lodash references remaining (lower is better tree-shaking); M8 = JS chunk count; M9 = sourcemap/JS ratio.}\label{tab:quality}
\begin{tabular}{@{}lrrrrrrrr@{}}
\toprule
 & \multicolumn{2}{c}{\textbf{M5 raw (KB)}} & \multicolumn{2}{c}{\textbf{M6 gzip (KB)}} & \multicolumn{2}{c}{\textbf{M7 tree-shake}} & \multicolumn{2}{c}{\textbf{M8 chunks}} \\
\cmidrule(lr){2-3}\cmidrule(lr){4-5}\cmidrule(lr){6-7}\cmidrule(lr){8-9}
\textbf{Tool} & \textbf{m-500} & \textbf{xl-5K} & \textbf{m-500} & \textbf{xl-5K} & \textbf{m-500} & \textbf{xl-5K} & \textbf{m-500} & \textbf{xl-5K} \\
\midrule
esbuild & 6{,}378 & 163{,}396 & 395 & 1{,}818 &  6 &  6 & 31 & 119 \\
Vite    & 1{,}611 &   8{,}843 & 320 & 1{,}064 &  3 &  3 & 19 &  84 \\
Rspack  & 3{,}377 &  27{,}952 & 504 & 3{,}176 &  2 &  2 &  5 &  12 \\
Rollup  & 4{,}460 &  25{,}327 & 612 & 2{,}553 & 17 & 17 &  6 &  13 \\
Webpack & 2{,}496 &  18{,}874 & 362 & 1{,}786 &  1 &  1 &  6 &  13 \\
\bottomrule
\end{tabular}
\end{table}

\textbf{Bundle size (M5/M6).}
Vite produces the smallest gzipped output at every scale: 320\,KB at m-500 versus 395\,KB for esbuild (+23\%) and 612\,KB for Rollup (+91\%).
At xl-5{,}000, Vite's gzipped output is 1{,}064\,KB---1.7$\times$ smaller than esbuild and 3.0$\times$ smaller than Rspack.
The large raw-size difference for esbuild (163\,MB at xl-5K) suggests less aggressive minification compared to the Rust/JS tools.

\textbf{Tree-shaking (M7).}
Webpack achieves the best tree-shaking with only 1 residual lodash reference across all scales.
Rspack (2 references) and Vite (3 references) follow closely.
esbuild retains 5--6 references, and Rollup retains 17---an unexpected result for a tool historically known for pioneering tree-shaking, likely due to its conservative handling of side effects in lodash-es re-exports.

\textbf{Code splitting (M8) and sourcemaps (M9).}
esbuild produces the most granular splitting (119 chunks at xl-5K), followed by Vite (84).
Rspack, Rollup, and Webpack converge at 12--13 chunks at xl-5K.
M9 (sourcemap coverage) shows a 1:1 ratio of sourcemaps to JS files for all tools at all scales---all tools generate complete sourcemaps.

\medskip\noindent\textbf{Finding 2.} \textit{Output quality differences are practically smaller than build speed differences. Vite produces the smallest bundles; Webpack achieves the best tree-shaking. All tools generate complete sourcemaps. The choice of bundler primarily affects build speed and memory, not output quality.}


\subsection{RQ3: Resource Consumption and Scaling}\label{sec:rq3}

\subsubsection{M10: Peak Memory (RSS)}

Table~\ref{tab:resources} presents memory and CPU results.

\begin{table}[htbp]
\centering
\caption{Resource consumption: M10 (peak RSS in MB) and M11 (CPU time in seconds), median of 10 runs. Bold indicates the lowest (best) value at each scale.}\label{tab:resources}
\begin{tabular}{@{}lrrrrr|rrrrr@{}}
\toprule
 & \multicolumn{5}{c|}{\textbf{M10: Peak Memory (MB)}} & \multicolumn{5}{c}{\textbf{M11: CPU Time (s)}} \\
\cmidrule(lr){2-6}\cmidrule(lr){7-11}
\textbf{Tool} & \textbf{xs} & \textbf{s} & \textbf{m} & \textbf{l} & \textbf{xl} & \textbf{xs} & \textbf{s} & \textbf{m} & \textbf{l} & \textbf{xl} \\
\midrule
esbuild & \textbf{101} & \textbf{141} & \textbf{124} & \textbf{177} & \textbf{287} & \textbf{1.41} & \textbf{1.75} & 3.43 & 8.80 & 28.67 \\
Vite    & 234 & 387 & 551 & 924 & 1{,}687 & 1.95 & 2.62 & \textbf{3.29} & \textbf{6.06} & \textbf{11.10} \\
Rspack  & 193 & 279 & 506 & 1{,}120 & 2{,}097 & 2.15 & 3.00 & 4.78 & 11.64 & 25.47 \\
Rollup  & 408 & 602 & 1{,}088 & 2{,}076 & 3{,}072 & 5.70 & 7.95 & 14.98 & 432.03 & 1{,}736.72 \\
Webpack & 479 & 897 & 2{,}124 & 4{,}832 & 5{,}925 & 7.81 & 10.33 & 20.89 & 76.99 & 214.52 \\
\bottomrule
\end{tabular}
\end{table}

esbuild uses the least memory at every scale, peaking at just 287\,MB for 5{,}000 modules---5.9$\times$ less than Vite (1{,}687\,MB) and 20.7$\times$ less than Webpack (5{,}925\,MB).
This reflects Go's efficient memory model with value types and deterministic garbage collection.

A notable anomaly: esbuild's memory at s-200 (141\,MB) is \emph{higher} than at m-500 (124\,MB).
This is consistent across all 10 runs (tight standard deviation) and likely reflects Go's GC behavior, where a smaller-but-more-fragmented heap may report higher RSS than a larger-but-contiguous allocation at a slightly bigger scale.
We report this as a genuine, reproducible finding.

Webpack's memory grows most aggressively: from 479\,MB (xs-50) to 5{,}925\,MB (xl-5{,}000), a 12.4$\times$ increase, making it impractical for CI environments with typical 4--8\,GB memory limits at enterprise project scales.

\subsubsection{M11: CPU Time}

esbuild is the most CPU-efficient at small scales ($\leq$200 modules): 1.41\,s at xs-50 and 1.75\,s at s-200.
The esbuild--Vite difference at these scales is not statistically significant ($p_\text{adj} = 0.052$ at xs-50, $p_\text{adj} = 0.284$ at s-200).
At $\geq$500 modules, Vite becomes the most CPU-efficient: 3.29\,s at m-500 ($<$ esbuild's 3.43\,s) and 11.10\,s at xl-5{,}000 ($<$ esbuild's 28.67\,s)---mirroring the M2 crossover.

Rollup exhibits a dramatic CPU explosion at large scales: 432\,s at l-2{,}000 and 1{,}737\,s (29 minutes) at xl-5{,}000.
This is 156$\times$ Vite's CPU time at xl-5{,}000 and is attributable to Rollup's exhaustive Temporal Dead Zone (TDZ) analysis in its JavaScript-based scope-hoisting pass.

\subsubsection{M12: Scaling Regression}

Table~\ref{tab:scaling} presents the scaling regression results.

\begin{table}[htbp]
\centering
\caption{Scaling regression: linear vs.\ log-linear fit of M2 medians across five Tier~1 scales. All tools are best fit by a linear model.}\label{tab:scaling}
\begin{tabular}{@{}lccrc@{}}
\toprule
\textbf{Tool} & $R^2_\text{linear}$ & $R^2_\text{log-linear}$ & \textbf{Slope (ms / 1{,}000 mod.)} & \textbf{Best Fit} \\
\midrule
Vite     & 0.997 & 0.776 &    954 & linear \\
Rspack   & 1.000 & 0.769 &  2{,}553 & linear \\
esbuild  & 0.972 & 0.637 &  2{,}888 & linear \\
Rollup   & 0.973 & 0.641 & 19{,}433 & linear \\
Webpack  & 0.995 & 0.711 & 22{,}770 & linear \\
\bottomrule
\end{tabular}
\end{table}

\textbf{All five tools scale linearly} with module count ($R^2 \geq 0.97$), not logarithmically.
This is a significant finding: linear scaling means that build times at 10{,}000 or 20{,}000 modules can be reliably extrapolated from our measurements.

The slopes vary by 23.9$\times$: Vite adds only 954\,ms per 1{,}000 additional modules, compared to 22{,}770\,ms for Webpack.
Rspack (2{,}553\,ms) and esbuild (2{,}888\,ms) have similar slopes, while Rollup (19{,}433\,ms) is close to Webpack, confirming that implementation language (Rust/Go vs.\ JavaScript) is the primary determinant of scaling efficiency.

\medskip\noindent\textbf{Finding 3.} \textit{Build time scales linearly for all tools ($R^2 \geq 0.97$). Vite's scaling slope is 23.9$\times$ lower than Webpack's. esbuild uses 5.9--20.7$\times$ less memory than competitors. At enterprise scale ($\geq$5{,}000 modules), Rollup's CPU consumption becomes prohibitive (29 minutes).}


\subsection{Tier~2: External Validation}\label{sec:tier2}

Table~\ref{tab:tier2} compares tool rankings between Tier~1 (s-200, the closest scale to Bulletproof React's 102 components) and Tier~2 (bp-react).

\begin{table}[htbp]
\centering
\caption{Tier~1 (s-200) vs.\ Tier~2 (Bulletproof React) ranking comparison. Arrows show rank order from fastest/lowest to slowest/highest. Rankings match exactly for M2 and M10; M11 has one adjacent-rank swap ($\Delta = 0.27$\,s).}\label{tab:tier2}
\begin{tabular}{@{}llll@{}}
\toprule
\textbf{Metric} & \textbf{Tier~1 (s-200)} & \textbf{Tier~2 (bp-react)} & \textbf{Match?} \\
\midrule
M2 (build time)  & eb $<$ vi $<$ rs $<$ ro $<$ wp & eb $<$ vi $<$ rs $<$ ro $<$ wp & \checkmark \\
M10 (peak RSS)   & eb $<$ rs $<$ vi $<$ ro $<$ wp & eb $<$ rs $<$ vi $<$ ro $<$ wp & \checkmark \\
M11 (CPU time)   & eb $<$ vi $<$ rs $<$ ro $<$ wp & eb $<$ rs $<$ vi $<$ ro $<$ wp & $\sim$\footnote{Rspack and Vite swap ranks; median gap = 0.27\,s (4.21\,s vs.\ 4.48\,s). Cliff's $\delta = 0.10$ (negligible).} \\
\bottomrule
\multicolumn{4}{@{}l}{\footnotesize eb = esbuild, vi = Vite, rs = Rspack, ro = Rollup, wp = Webpack.}
\end{tabular}
\end{table}

Rankings match exactly for M2 (production build time) and M10 (peak memory).
For M11 (CPU time), Rspack and Vite swap adjacent ranks, but their Tier~2 medians differ by only 0.27\,s (4.21\,s vs.\ 4.48\,s), with a negligible Cliff's delta of 0.10.
This swap is well within measurement noise and does not contradict the Tier~1 findings.

The Tier~1/Tier~2 agreement provides evidence that our calibrated synthetic projects produce rankings that generalize to a real-world application with production dependencies (React Query, Zustand, Radix UI, Tailwind CSS).

\medskip\noindent\textbf{Finding 4.} \textit{Tier~2 (real-world OSS) rankings match Tier~1 (synthetic) rankings for M2 and M10 exactly, with only one negligible adjacent-rank swap for M11. The two-tier strategy validates external generalizability.}

%% ═══════════════════════════════════════════════════════════════
%% SECTION 6: DISCUSSION (placeholder)
%% ═══════════════════════════════════════════════════════════════
\section{Discussion}\label{sec:discussion}

\textit{[To be written --- see WRITING-PLAN.md \S6]}

%% ═══════════════════════════════════════════════════════════════
%% SECTION 7: THREATS TO VALIDITY (placeholder)
%% ═══════════════════════════════════════════════════════════════
\section{Threats to Validity}\label{sec:threats}

\textit{[To be written --- see WRITING-PLAN.md \S7]}

%% ═══════════════════════════════════════════════════════════════
%% SECTION 8: CONCLUSION (placeholder)
%% ═══════════════════════════════════════════════════════════════
\section{Conclusion and Future Work}\label{sec:conclusion}

\textit{[To be written --- see WRITING-PLAN.md \S8]}


%% ─── Data Availability ───
\section*{Data Availability}
The complete replication package---including the project generator, bundler configurations, measurement scripts, raw data (277 CSV files, 1{,}534 data points), analysis code, and all results---is publicly available at [Zenodo DOI to be assigned upon submission].

%% ─── Declaration of Generative AI Use ───
\section*{Declaration of Generative AI and AI-Assisted Technologies}
During the preparation of this manuscript, the authors used AI-assisted tools for code generation (measurement scripts) and statistical analysis verification.
The authors reviewed and validated all AI-generated content and take full responsibility for the accuracy and integrity of the published work.

%% ─── Declaration of Competing Interests ───
\section*{Declaration of Competing Interests}
The authors declare that they have no known competing financial interests or personal relationships that could have appeared to influence the work reported in this paper.
This study was conducted independently with no funding or affiliation with any of the bundler projects evaluated.

\bibliographystyle{elsarticle-num}
\bibliography{references}

\end{document}
